\documentclass[11pt,a4paper]{article}

% ─── Packages ───────────────────────────────────────────────
\usepackage[utf8]{inputenc}
\usepackage[T1]{fontenc}
\usepackage{lmodern}
\usepackage[margin=2.3cm]{geometry}
\usepackage{amsmath,amssymb}
\usepackage{graphicx}
\usepackage{xcolor}
\usepackage{hyperref}
\usepackage{booktabs}
\usepackage{array}
\usepackage{tabularx}
\usepackage{enumitem}
\usepackage{fancyhdr}
\usepackage{titlesec}
\usepackage{listings}
\usepackage{tcolorbox}
\usepackage{fontawesome5}
\usepackage{parskip}

% ─── Colours (Earth Tones palette of the app) ───────────────
\definecolor{accent}{HTML}{2E7D32}       % green (app accent)
\definecolor{brown}{HTML}{5D4037}        % dark brown (toolbar)
\definecolor{brownmid}{HTML}{6B4C3A}     % body background
\definecolor{brownlight}{HTML}{BCAAA4}   % panel
\definecolor{cream}{HTML}{EFEBE9}        % surface
\definecolor{paper}{HTML}{FFF8F0}        % plot background
\definecolor{warnhl}{HTML}{9A6700}
\definecolor{redhl}{HTML}{8B2E2E}
\definecolor{border}{HTML}{BCAAA4}

% ─── Hyperref ───────────────────────────────────────────────
\hypersetup{
  colorlinks=true,
  linkcolor=accent,
  urlcolor=brown,
  citecolor=accent,
  pdftitle={Conditional Dimensional Explorer - User Manual},
  pdfsubject={User Manual}
}

% ─── Header / Footer ────────────────────────────────────────
\pagestyle{fancy}
\fancyhf{}
\fancyhead[L]{\small\textcolor{brown}{\textbf{Conditional Dimensional Explorer}}}
\fancyhead[R]{\small Earth Tones Edition}
\fancyfoot[C]{\small\thepage}
\renewcommand{\headrulewidth}{0.4pt}

% ─── Section styling ────────────────────────────────────────
\titleformat{\section}
  {\Large\bfseries\color{brown}}{\thesection}{0.6em}{}
\titleformat{\subsection}
  {\large\bfseries\color{black!80}}{\thesubsection}{0.5em}{}
\titleformat{\subsubsection}
  {\normalsize\bfseries\color{black!70}}{\thesubsubsection}{0.4em}{}

% ─── Custom boxes ───────────────────────────────────────────
\newtcolorbox{notebox}[1][]{
  colback=cream, colframe=brown, boxrule=0.6pt,
  arc=2mm, left=4mm, right=4mm, top=2mm, bottom=2mm,
  fonttitle=\bfseries\small, title={#1}
}
\newtcolorbox{tipbox}[1][]{
  colback=accent!6, colframe=accent, boxrule=0.6pt,
  arc=2mm, left=4mm, right=4mm, top=2mm, bottom=2mm,
  fonttitle=\bfseries\small, title={#1}
}
\newtcolorbox{warnbox}[1][]{
  colback=orange!6, colframe=warnhl, boxrule=0.6pt,
  arc=2mm, left=4mm, right=4mm, top=2mm, bottom=2mm,
  fonttitle=\bfseries\small, title={#1}
}

% ─── Listings ───────────────────────────────────────────────
\lstset{
  basicstyle=\ttfamily\small,
  backgroundcolor=\color{cream},
  frame=single,
  rulecolor=\color{border},
  breaklines=true,
  columns=fullflexible,
  keepspaces=true,
}

\setlength{\parindent}{0pt}
\setlength{\parskip}{5pt}

\begin{document}

% ════════════════════════════════════════════════════════════
% TITLE PAGE
% ════════════════════════════════════════════════════════════
\begin{titlepage}
\centering
\vspace*{3cm}

{\Huge\bfseries\color{accent} Conditional Dimensional Explorer}\\[0.4cm]
{\LARGE\bfseries User Manual}\\[1.5cm]

{\large Earth Tones Edition}\\[0.3cm]
{\large Interactive Conditional Data Visualization Tool}\\[3cm]

\begin{tabular}{rl}
\textbf{Application type:} & Single-page web application (HTML) \\
\textbf{Plotting engine:} & Plotly.js 3.1.0 \\
\textbf{Views:} & Histogram, Scatter (2D), Scatter3D \\
\textbf{Conditioning:} & Slider-based, distance-threshold filter \\
\textbf{Data input:} & CSV upload, synthetic generators \\
\textbf{Exports:} & CSV \\
\end{tabular}

\vfill

{\small This document explains how to operate the Conditional Dimensional Explorer: how to load or generate data, condition the point cloud, switch between one-, two- and three-dimensional views, and export results.}

\end{titlepage}

% ════════════════════════════════════════════════════════════
\tableofcontents
\newpage

% ════════════════════════════════════════════════════════════
\section{Introduction}
% ════════════════════════════════════════════════════════════

The \textbf{Conditional Dimensional Explorer} is a browser-based tool for the interactive, conditional visualization of multivariate data. It transforms every variable into a uniform [0,\,1] scale (empirical CDF) and then lets you explore subsets of the data selected by \emph{conditioning} variables. The points closest to a chosen conditioning point are highlighted, while the rest fade toward transparency.

The application supports three visualization modes depending on how many variables you choose to represent:

\begin{itemize}[leftmargin=1.4em]
  \item \textbf{1 variable} --- a histogram of the selected variable's original values.
  \item \textbf{2 variables} --- a scatter plot in the unit square.
  \item \textbf{3 variables} --- an interactive 3D scatter plot.
\end{itemize}

\subsection{Key features}

\begin{itemize}[leftmargin=1.4em]
  \item Load any numeric CSV file (comma, semicolon or tab separated).
  \item Generate two synthetic datasets: a fixed 15-variable structured sample and a 10-variable realistic correlated sample.
  \item Condition on any subset of variables using sliders.
  \item Fade-out by distance: the ``distance threshold'' slider controls how sharply the cloud fades.
  \item Save the active dataset as CSV.
  \item View the data as a table and inspect a Pearson correlation matrix.
  \item Full-screen mode and an exit button.
\end{itemize}

\begin{notebox}[No installation required]
The application is a single HTML file. It requires an internet connection only on first load, because it fetches Plotly.js from a CDN. All computation runs locally in your browser.
\end{notebox}

% ════════════════════════════════════════════════════════════
\section{How the Visualization Works}
% ════════════════════════════════════════════════════════════

\subsection{Empirical CDF normalization}

Each variable is transformed to a uniform scale on $[0,1]$ using its \textbf{empirical cumulative distribution function} (ECDF):
\begin{equation}
  q_{ij} = \frac{\operatorname{rank}(x_{ij})}{n+1}
\end{equation}
where $n$ is the number of rows and $\operatorname{rank}(x_{ij})$ is the rank of the $j$-th column value in row $i$ (with rank 1 the smallest). This makes every axis comparable regardless of units or scale.

\subsection{Conditioning and the distance threshold}

If you select $k$ conditioning variables $\{c_1,\dots,c_k\}$, you also get $k$ sliders in the right panel. For each point, a distance in the normalized space is computed:

\begin{equation}
  d_i = \sqrt{\sum_{m=1}^{k} \left(q_{i c_m} - \theta_{c_m}\right)^2}
\end{equation}
where $\theta_{c_m}$ is the slider value of the conditioning variable. Points with $d_i < \tau$ (with $\tau$ the \emph{distance threshold}) are shown with an opacity

\begin{equation}
  \alpha_i = \max\left(0.2,\; 1 - \frac{d_i}{\tau}\right).
\end{equation}

Points with $d_i \ge \tau$ are hidden. Lowering $\tau$ focuses on a small neighbourhood; raising it shows more points.

\subsection{The three view modes}

\begin{center}
\begin{tabularx}{\textwidth}{lX}
\toprule
\textbf{Represented} & \textbf{What you see} \\
\midrule
1 variable & A histogram of that variable's raw (unnormalized) values, filtered by the current conditioning. \\
2 variables & A 2D scatter plot of their ECDF coordinates $[0,1]\times[0,1]$, with a green frame. \\
3 variables & A 3D scatter plot of their ECDF coordinates inside a unit cube. \\
\bottomrule
\end{tabularx}
\end{center}

\begin{tipbox}[The axes are always ECDF coordinates in 2D/3D]
In the 2D and 3D views, the axes show \textbf{normalized} values between 0 and 1, so point positions on every axis are comparable. In the histogram view (1 variable), the raw original values are shown on the x-axis.
\end{tipbox}

% ════════════════════════════════════════════════════════════
\section{Getting Started}
% ════════════════════════════════════════════════════════════

\subsection{Opening the application}

Open the HTML file in a modern browser (Chrome, Firefox, Edge, Safari). On first load the app automatically generates the fixed 15-variable structured sample, so you can start exploring immediately.

\subsection{Interface overview}

\begin{center}
\begin{tabularx}{\textwidth}{lX}
\toprule
\textbf{Region} & \textbf{Content} \\
\midrule
Top toolbar & Title, file upload, sample-size, data generators, variable selection, save, table, correlation, fullscreen and exit buttons. \\
Central plot & The active figure (histogram, scatter or 3D). \\
Right panel & Distance threshold, point counter, and conditioning sliders. \\
Modals & Variable selection, data table, correlation matrix. \\
\bottomrule
\end{tabularx}
\end{center}

\subsection{The point counter}

Below the distance-threshold slider you see a badge of the form \texttt{2000 / 2000 pts}. The first number is how many points currently satisfy the distance threshold; the second is the total number of rows. This gives immediate feedback on how restrictive your conditioning is.

% ════════════════════════════════════════════════════════════
\section{The Toolbar in Detail}
% ════════════════════════════════════════════════════════════

\begin{center}
\begin{tabularx}{\textwidth}{lX}
\toprule
\textbf{Button} & \textbf{Action} \\
\midrule
\faFolderOpen{} \textbf{Upload CSV} & Opens a file picker. Loads a numeric CSV (auto-detects comma, semicolon or tab). \\
\textbf{n =} (field) & Sample size for the synthetic generators (200--20\,000). \\
\faBolt{} \textbf{Data Initialize} & Generates the fixed 15-variable structured dataset. \\
\faDna{} \textbf{Realistic \textpm\ corr (10 vars)} & Generates a 10-variable, correlated sample (Income, Education, Debt, \dots). \\
\faSlidersH{} \textbf{Variable selection} & Opens the modal to choose represented and conditioning variables. \\
\faSave{} \textbf{Save CSV} & Downloads the current dataset. \\
\faClipboard{} \textbf{Show data table} & Opens a modal with the first 500 rows. \\
\faChartLine{} \textbf{Correlation matrix} & Opens a Pearson correlation heatmap. \\
\faExpand{} \textbf{Fullscreen} & Toggles fullscreen. \\
\faPowerOff{} \textbf{Exit} & Attempts to close the tab (subject to browser rules). \\
\bottomrule
\end{tabularx}
\end{center}

\begin{notebox}[Sample size only affects the generators]
The \texttt{n =} field controls how many rows the two synthetic generators produce. It has no effect on a CSV you upload; the field updates to reflect the loaded row count.
\end{notebox}

% ════════════════════════════════════════════════════════════
\section{Loading Data}
% ════════════════════════════════════════════════════════════

\subsection{From a CSV file}

Click \textbf{Upload CSV} and choose a file. The app:

\begin{enumerate}[leftmargin=1.4em]
  \item Detects the delimiter (comma, semicolon or tab) from the first non-empty line.
  \item Parses the file (supports quoted fields and quoted line breaks).
  \item Keeps only rows whose cells are all numeric and complete.
  \item If some rows are invalid, shows a confirmation dialog with the first affected row numbers before loading the valid subset.
\end{enumerate}

The first row must be the header (variable names); every subsequent row must have exactly as many numeric values as the header.

\subsubsection*{Decimal comma}

When the file is separated by semicolon or tab, decimal commas are understood (e.g., \texttt{3,14}). For comma-separated files the decimal separator must be a point.

\subsection{Replacing the active dataset}

Whenever you load or generate a new dataset, if a dataset is already active the app asks you whether you want to save it as CSV first. This prevents accidental loss of work.

\begin{warnbox}[Irreversible replacement]
If you choose not to save, the active dataset, its plots, its modals and the conditioning state are discarded. The new data is installed as a fresh state: two represented variables, one conditioning variable, threshold 0.80 and all sliders at 0.50.
\end{warnbox}

\subsection{Synthetic datasets}

Two built-in generators are available:

\begin{itemize}[leftmargin=1.4em]
  \item \textbf{Data Initialize} --- a fixed 15-variable structure named after academic subjects (math, physics, chemistry, \dots, coordination), generated from a five-factor latent model with a fixed loading matrix and a small correlated-factor matrix.
  \item \textbf{Realistic \textpm\ corr} --- 10 variables (Income, Education, Debt, Savings, Stress, Health, SocialScore, RiskTolerance, WorkHours, LeisureTime) generated from a 10$\times$10 target correlation matrix, adjusted to positive-definite via a small ridge, then drawn through a Cholesky factor.
\end{itemize}

Both generators clamp every value to $[0, 100]$ and round to one decimal.

% ════════════════════════════════════════════════════════════
\section{Choosing Variables}
% ════════════════════════════════════════════════════════════

Click \textbf{Variable selection} to open the modal. It has two lists:

\begin{itemize}[leftmargin=1.4em]
  \item \textbf{Represented variables} --- one, two or three. These are what you actually plot.
  \item \textbf{Conditioning variables} --- any number. These become the sliders.
\end{itemize}

\subsection{Rules}

\begin{itemize}[leftmargin=1.4em]
  \item You must select at least 1 represented variable.
  \item You may select at most 3 represented variables.
  \item Conditioning variables cannot overlap with represented variables; the app disables those entries in the conditioning list.
  \item You can choose \textbf{None (unconditional)} to remove all conditioning sliders.
\end{itemize}

Click \textbf{Apply \& close} to accept. The app rebuilds the sliders and the plot immediately, and then requests fullscreen (a convenience feature that some browsers may block if no user gesture is detected).

\begin{tipbox}[Typical workflows]
\textbf{One-dimensional}: represent 1 variable and condition on 2--3 others to see how its distribution shifts. \\
\textbf{Two-dimensional}: represent 2 variables and condition on 1--2 others to see the shape of the cloud from nearby slices. \\
\textbf{Three-dimensional}: represent 3 variables, no conditioning, to inspect the full cube.
\end{tipbox}

% ════════════════════════════════════════════════════════════
\section{Conditioning and the Threshold}
% ════════════════════════════════════════════════════════════

Once you have at least one conditioning variable, its slider appears in the right panel. Each slider ranges from 0 to 1 (ECDF units).

\subsection{The distance threshold}

The top slider of the right panel, labeled \emph{distance threshold}, ranges from 0 to 3. Its value $\tau$ determines:

\begin{itemize}[leftmargin=1.4em]
  \item \textbf{Hidden points}: any point with $d_i \ge \tau$.
  \item \textbf{Opacity}: $\alpha_i = \max(0.2,\, 1 - d_i/\tau)$ for the points kept.
\end{itemize}

Consequently, \textbf{smaller} $\tau$ means a tighter focus (fewer, brighter points near the conditioning point) and \textbf{larger} $\tau$ means a looser focus (more points, more uniform alpha).

\subsection{The point counter}

The badge under the threshold shows \texttt{kept / total pts}. Use it to quantify how restrictive the current setting is.

\begin{tipbox}[Getting a useful focus]
Start with the threshold at 0.80. If the badge shows too many points (e.g., more than half the dataset), lower it; if it shows too few, raise it. When you move a conditioning slider, the badge updates instantly.
\end{tipbox}

% ════════════════════════════════════════════════════════════
\section{The Three View Modes}
% ════════════════════════════════════════════════════════════

\subsection{Histogram (1 represented variable)}

The histogram counts the raw, unnormalized values of that variable among the points that pass the conditioning filter. Its x-axis shows the original values, not the ECDF coordinates. This is the fastest way to see how a variable's distribution responds to conditioning.

\subsection{Scatter plot (2 represented variables)}

Points are shown inside the unit square $[0,1]\times[0,1]$ in ECDF coordinates, with a green frame indicating the boundaries. The transparency of each point reflects how close it is to the conditioning point; the axis limits are always $[0,1]$.

\subsection{3D scatter (3 represented variables)}

An interactive 3D scatter of the ECDF coordinates inside the unit cube. The Plotly camera controls are enabled (drag to rotate, scroll to zoom). The camera position is preserved while you move the sliders, so you can focus on the region of interest without the view jumping on every update.

\begin{notebox}[Appearance is fixed]
In the 2D and 3D views, the axes are always set to $[0,1]$ so that comparisons across slices remain meaningful. Resetting the view is done through the Plotly modebar's ``Home'' button.
\end{notebox}

% ════════════════════════════════════════════════════════════
\section{The Data Table and the Correlation Matrix}
% ════════════════════════════════════════════════════════════

\subsection{Data table}

Click \textbf{Show data table} to see the first 500 rows. Columns are exactly the header names, and numeric values are rounded to three decimals for readability. If the dataset has more than 500 rows, a footer row states how many are hidden.

\subsection{Correlation matrix}

Click \textbf{Correlation matrix} to compute a Pearson correlation for every pair of variables and render it as a heatmap. The scale is $[-1,1]$ (blue to red), with the numerical value printed in each cell. This is useful to check whether a conditioning variable is (or is not) correlated with a represented one before selecting it.

% ════════════════════════════════════════════════════════════
\section{Saving Your Work}
% ════════════════════════════════════════════════════════════

Click \textbf{Save CSV} to download the current dataset as a comma-separated file. All variable names are preserved and each value is written verbatim (so re-importing gives the same data). A UTF-8 byte-order mark is prepended so that Excel recognizes the encoding.

\begin{warnbox}[Save before generating]
Loading a CSV or clicking a generator will prompt you to save the active data if it exists. Always accept the prompt if you want to keep your work.
\end{warnbox}

% ════════════════════════════════════════════════════════════
\section{Fullscreen and Exit}
% ════════════════════════════════════════════════════════════

\textbf{Fullscreen} toggles the browser's fullscreen mode on the whole page. It is also triggered automatically when you apply a new variable selection.

\textbf{Exit} asks for confirmation and then tries to close the tab. Because of browser security, a page can only close itself if it was opened by a script. If it cannot close, an alert reminds you to close the tab manually.

% ════════════════════════════════════════════════════════════
\section{Troubleshooting}
% ════════════════════════════════════════════════════════════

\begin{description}[leftmargin=0em, style=nextline]
  \item[``Plotly library could not be loaded'']
  The page cannot reach the Plotly CDN. Connect to the internet and reload the page. The application cannot plot without Plotly.

  \item[``The file must contain at least two variables'']
  Your CSV has only one column. Make sure the header row contains at least two names.

  \item[``No valid numeric rows found'']
  One or more rows are not fully numeric. If decimal commas are used, make sure the file is semicolon- or tab-separated (not comma-separated).

  \item[``Unbalanced quotes'']
  A quoted field is missing its closing quote. Fix the CSV and try again.

  \item[Only a few points are visible]
  The distance threshold is too small. Raise it, or move the conditioning sliders toward the center (0.5).

  \item[The selected variable does not appear as an axis]
  It is probably in the conditioning list. The app disables overlapping entries in the modal; re-open the modal and adjust the selection.

  \item[The full screen does not open automatically]
  Some browsers block automatic fullscreen if it is triggered outside a user gesture. Use the Fullscreen button instead.
\end{description}

% ════════════════════════════════════════════════════════════
\section{Quick Reference}
% ════════════════════════════════════════════════════════════

\begin{center}
\begin{tabularx}{\textwidth}{lX}
\toprule
\textbf{Task} & \textbf{Action} \\
\midrule
Load a CSV & Upload CSV \textrightarrow{} pick file. \\
Generate default data & Data Initialize. \\
Generate correlated data & Realistic \textpm\ corr (10 vars). \\
Change plotted variables & Variable selection \textrightarrow{} pick 1, 2 or 3 \textrightarrow{} Apply. \\
Change conditioning & Same modal; pick conditioning variables. \\
Focus tighter & Lower the distance-threshold slider. \\
Inspect rows & Show data table. \\
Inspect correlations & Correlation matrix. \\
Export data & Save CSV. \\
Fullscreen & Fullscreen button (or automatically after Apply). \\
Exit & Exit button. \\
\bottomrule
\end{tabularx}
\end{center}

\begin{notebox}[Recommended workflow]
\begin{enumerate}
  \item Load or generate the dataset.
  \item Inspect the correlation matrix to understand relationships.
  \item Open Variable selection; choose 2 represented variables and 1 conditioning variable as a starting point.
  \item Move the conditioning slider to explore different slices.
  \item Lower the distance threshold to focus on a local neighbourhood.
  \item Switch to 3 represented variables for a full cube view.
  \item Save the CSV if you have edited the data.
\end{enumerate}
\end{notebox}

% ════════════════════════════════════════════════════════════
\section*{Version Notes}
\addcontentsline{toc}{section}{Version Notes}
% ════════════════════════════════════════════════════════════

This manual refers to the version stored as \texttt{ConditionalExplorerFINALNEW.html}. Its distinguishing features relative to earlier drafts include:

\begin{itemize}[leftmargin=1.4em]
  \item Full replacement semantics when loading a new dataset (old state, plots and modals are cleared).
  \item Safe CSV replacement with an optional save-before-replace prompt.
  \item Robust CSV parser with delimiter auto-detection and per-row validation.
  \item Plotly camera preservation across slider updates in the 3D view.
  \item ``None (unconditional)'' option in the conditioning list.
  \item Automatic fullscreen when the user applies a new variable selection.
  \item Distance-threshold semantics: points with $d_i \ge \tau$ are hidden, opacity decays as $1 - d_i/\tau$.
\end{itemize}

\end{document}